\par\Needspace{7\baselineskip}\bigskip
\begin{center}
\SourceLeaf{442}
\pdfbookmark[2]{Gauss to Hansen, Göttingen, 11 March 1843}{pallas-brief-18430311}
\textsc{Gauss} to \textsc{Hansen}, Göttingen, 11 March 1843.
\end{center}
\noindent\hspace*{\parindent}From one week to the next I have put off thanking you for kindly sending me the Berliner Monatsbericht (which came to my attention through the usual channels at about the same time), because I hoped to find some time to go a little further into the subject in my reply. Unfortunately that hope has been disappointed, and even in the coming holidays, for which so much unfinished business and so many new demands have already accumulated, I can scarcely hope for any leisure.
\par I shall therefore confine myself to expressing my pleasure that in your perturbation calculations you proceed in much the same way as I did more than 30 years ago in my extensive calculations of the perturbations of Pallas, in that you dispense altogether with the use of series in the eccentricities and inclinations. Of course, for the perturbations of comets you also needed quite distinct investigations for which there was no occasion in the case of Pallas, and on the whole you took different paths in many other respects. For the perturbations of Pallas due to Jupiter I deliberately used the method of variable elements. Although this entails treating 6 elements, whereas the other procedure has only half as many, I preferred the former for practical use. The calculation need be made only once for each year (for a single day). For this purpose I constructed a special auxiliary table with which the calculation can be completed in a very short time by comparison, although in all 801 equations were taken into account (1602 if the sine and cosine terms of the same argument are counted separately. Every one whose coefficient exceeds \(0''{,}1\) is included)\footnote{It is probably unnecessary to remark that the work of which I am now speaking is \textls[180]{entirely} different from the calculation by quadrature which I also carried through for the entire period from 1802 to about 1818 or 1820. I call this the special calculation and the other the general calculation, and the latter found in the former a check of the utmost necessity for calculations of such scope.}. Very simple means then suffice to make the result adequate for the whole \SourceLeaf{443}observing season without any sacrifice of precision. The perturbation due to Saturn was also calculated, and that due to Mars was begun by a substantially different method, but was left unfinished. Other occupations, which kept branching out, later left me no opportunity at all to return to that work. It remains to be seen whether at my age I shall have the leisure and inclination to immerse myself in these matters again, since there are still so many other things that I would be just as reluctant to let go to waste. You are very fortunate to be in circumstances where you need not spread yourself thin, and it will in any event give me peace of mind that this branch of astronomy is in the best of hands with you. Whether I used anything in the perturbations due to Mars that was connected with the method of integration you indicated to me, I cannot now determine with certainty, having been away from these matters for nearly 25 years, but I am inclined to doubt it.\quad\mbox{. . . . .}
\par
